
\subsection{About Open Declaration}
(1897.)

STATE OF MINNESOTA; COUNTY OF HENNEPIN.
DISTRICT COURT; FOURTH JUDICIAL DISTRICT.

In the Matter of the Application of Nils C. Brun, Lauritz
F. Clausen, Lars Swenson, Orson A. Veblen and Ole
O. Aanstad, for leave to File an Information, in the
nature of quo warranto, against Sven Oftedal, Olaf
Hoff, Ole Paulson, Theodore Helgesen and Andrew
Knutson.

AFFIDAVIT OF SVEN OFTEDAL.

State of Minnesota, County of Hennepin—ss.

Sven Oftedal being first duly sworn, upon his oath
deposes and says: . . . .

At the time I thus became connected with Augsburg
Seminary there existed a difference of opinion among Nor-
wegian Lutherans respecting church polity and methods
of theological instruction.* Such difference of opinion has
ever since existed among them, and contributed largely
to the controversies which thereafter arose in the Confer-
ence, as will be shown in this affidavit. On the one hand
were the adherents of the “new school,” as it is called,
whose convictions were and are that the individual congre-
gation is “the church;”* and should act with other congrega-
tions for common purposes only, and reserve to itself con-
trol respecting questions pertaining to the administration
of its own affairs and those also bearing upon “purity of
doctrine,” as these words are understood among theologi-
ans. Those entertaining these views respecting church
polity call them “free church” ideas.* On the same side were
those who believed that ministers should be trained and
educated in accordance with these views,* and with refer-
ence to their future calling, throughout their collegiate or
preparatory course, as well as during their theological
course, as the same is commonly understood. On the other
side were the adherents of what is called the “old school,”
who opposed, to a greater or less extent, the views of the
new school, and desired, as it was believed, to partially
adopt or retain certain methods of church government*
prevailing in the established or state church of Norway.

The theological professors and instructors of Augsburg
Seminary, and its board of trustees now in control of the
seminary, have adhered to the new school as herein de-
fined, and have conducted the seminary in accordance with
its views. And in entering upon my work at Augsburg
Seminary I was actuated by a desire and willingness to
devote my entire life to the advancement of free church
ideas,* and what I then believed, and still believe, to be
the true purpose and end of a theological or divinity
school. As will be seen, these free church doctrines and
ideas thus entertained by the present theological profes-
sors of Augsburg Seminary, while antagonizing certain of
the ministers of the Conference, obtained for them the
earnest and enthusiastic support of a large majority of the
lay delegates participating in the annual meetings of the
society, and of the individual members of the various
congregations constituting that association. . . . .

In the year 1873, and ever since, there has been in this
country an organization of Norwegian Lutheran churches
officially entitled “The Synod of the Norwegian Evangeli-
cal Lutheran Church of America,” and commonly called
the Norwegian Synod, between which and the Conference
there existed in the year last mentioned, as well as in later
years, a somewhat bitter doctrinal controversy. The
doctrines and church polity of this Norwegian Synod has
always, as I believe, been extremely inimical to the growth
of a free church,* and on January 20th, 1874 Professor
Weenaas and this affiant caused to be published a writ-
ten document or statement signed by ourselves, the title
of which, translated into English, was “An Open Declara-
tion.” In this statement the church polity of the Norwe-
gian Synod was arraigned and criticised.* In this state-
ment were set forth the principles, already mentioned, by
which the subscribers intended to be guided in church
work, and by it they emphasized the necessity for the
maintenance and building up of Augsburg Seminary. This
document brought forth much adverse criticism from the
adherents of the said Norwegian Synod, and created some
dissatisfaction among a faction of the Conference itself.
This writing subsequently received official mention in the
proceedings of the Conference, and gave rise in the society
to opposition on the part of certain of its members to
Augsburg Seminary and its management. As will be seen,
this opposition, thus initiated, continued until the present
time, but was, from the year 1883, at least, until the for-
mation of what is called the United Church, at all times
in a hopeless minority; but after the union such minority,
through a combination with other element and churches
which were never a part of the Conference, caused to be
instituted and brought this application now pending. . . . .

SVEN OFTEDAL.

Subscribed and sworn to before me this 17th of March,
1897.

F. N. HENDRIX,
Notary Public, Hennepin County, Minn.

Notarial Seal